% GENERATED FILE. Edit canonical lessons, then run npm run generate:books.
% canonical-source-hash: fnv1a64:36a37fbc78c3b20c
% canonical-lessons: JA-C273-saiban, JA-C273-bakkin, JA-C273-hanketsu, JA-C273-uttaeru, JA-C273-shouko, JA-R273-review-268-273

\chapter{The Trial}
\label{ch:ja-the-trial}
\hlchaptermodality{273}

\begin{chapteropening}
\textbf{By the end of this chapter:} \emph{I can say what happened in court and what was paid.}

Complete the last lesson of chapter 273: I can say what happened in court and what was paid.
\end{chapteropening}

\section[\texorpdfstring{\ja{さいばん} (saiban)}{さいばん (saiban)}]{\ja{さいばん} (saiban) — a trial; a court case}

\label{lesson:JA-C273-saiban}



\begin{quote}
Before the new one: say the Japanese for to hide (something), then the Japanese for true.
\end{quote}

\subsection*{You'll want to know: \ja{さいばん}}
\textbf{\ja{さいばん}} — \emph{saiban} — \textquotedblleft{}a trial; a court case\textquotedblright{}.

\begin{quote}
\textbf{\ja{あした} \ja{さいばんが} \ja{あります。}} — \emph{ashita saiban ga arimasu} — \textquotedblleft{}The trial is tomorrow.\textquotedblright{}
\end{quote}

\begin{cousinweb}
\textbf{\ja{さいばん}} (\emph{saiban}) is Sino-Japanese (kango): two Chinese characters read with their on'yomi, readings borrowed from Chinese. They are \emph{sai}, \textquotedblleft{}judge\textquotedblright{} (Middle Chinese \emph{dzoj}) and \emph{ban}, \textquotedblleft{}judgement\textquotedblright{} (\emph{phanH}) (KANJIDIC2; Middle Chinese in Baxter's transcription). The second is voiced from \emph{han} to \emph{ban}.
\end{cousinweb}

\begin{grammarlens}[title={What you've built}]
One more word for telling what happened.
\end{grammarlens}

\subsection*{Your turn}
\begin{itemize}
\raggedright
  \item \emph{Say it:} \emph{saiban}
  \item \emph{Say it:} \emph{saiban}, once more
  \item \emph{Say it:} \emph{ashita saiban ga arimasu} — The trial is tomorrow.
  \item \emph{Recall:} say the Japanese for to hide (something), then the Japanese for true, then say \emph{saiban} again
\end{itemize}

\begin{tcolorbox}[breakable,colback=teal!4,colframe=teal!35!black,title={Before you move on}]
What does \ja{さいばん} mean? (\textquotedblleft{}A trial; a court case\textquotedblright{}.) Say it once more.
\end{tcolorbox}
\section[\texorpdfstring{\ja{ばっきん} (bakkin)}{ばっきん (bakkin)}]{\ja{ばっきん} (bakkin) — a fine; a penalty}

\label{lesson:JA-C273-bakkin}



\begin{quote}
Before the new one: say the Japanese for true, then the Japanese for a trial.
\end{quote}

\subsection*{You'll want to know: \ja{ばっきん}}
\textbf{\ja{ばっきん}} — \emph{bakkin} — \textquotedblleft{}a fine; a penalty\textquotedblright{}.

\begin{quote}
\textbf{\ja{ちちは} \ja{ばっきんを} \ja{はらいました。}} — \emph{chichi wa bakkin o haraimashita} — \textquotedblleft{}My father paid a fine.\textquotedblright{}
\end{quote}

\begin{cousinweb}
\textbf{\ja{ばっきん}} (\emph{bakkin}) is Sino-Japanese (kango): two Chinese characters read with their on'yomi, readings borrowed from Chinese. They are \emph{batsu}, \textquotedblleft{}penalty\textquotedblright{} (Middle Chinese \emph{bjot}) and \emph{kin}, \textquotedblleft{}gold\textquotedblright{} (\emph{kim}) (KANJIDIC2; Middle Chinese in Baxter's transcription). The first is shortened from \emph{batsu} to \emph{bak-} before \emph{kin}.
\end{cousinweb}

\begin{grammarlens}[title={What you've built}]
One more word for telling what happened.
\end{grammarlens}

\subsection*{Your turn}
\begin{itemize}
\raggedright
  \item \emph{Say it:} \emph{bakkin}
  \item \emph{Say it:} \emph{bakkin}, once more
  \item \emph{Say it:} \emph{chichi wa bakkin o haraimashita} — My father paid a fine.
  \item \emph{Recall:} say the Japanese for true, then the Japanese for a trial, then say \emph{bakkin} again
\end{itemize}

\begin{tcolorbox}[breakable,colback=teal!4,colframe=teal!35!black,title={Before you move on}]
What does \ja{ばっきん} mean? (\textquotedblleft{}A fine; a penalty\textquotedblright{}.) Say it once more.
\end{tcolorbox}
\section[\texorpdfstring{\ja{はんけつ} (hanketsu)}{はんけつ (hanketsu)}]{\ja{はんけつ} (hanketsu) — a verdict; a judgement}

\label{lesson:JA-C273-hanketsu}



\begin{quote}
Before the new one: say the Japanese for a trial, then the Japanese for a fine.
\end{quote}

\subsection*{You'll want to know: \ja{はんけつ}}
\textbf{\ja{はんけつ}} — \emph{hanketsu} — \textquotedblleft{}a verdict; a judgement\textquotedblright{}.

\begin{quote}
\textbf{\ja{はんけつは} \ja{らいしゅうです。}} — \emph{hanketsu wa raishū desu} — \textquotedblleft{}The verdict is next week.\textquotedblright{}
\end{quote}

\begin{cousinweb}
\textbf{\ja{はんけつ}} (\emph{hanketsu}) is Sino-Japanese (kango): two Chinese characters read with their on'yomi, readings borrowed from Chinese. They are \emph{han}, \textquotedblleft{}judgement\textquotedblright{} (Middle Chinese \emph{phanH}) and \emph{ketsu}, \textquotedblleft{}decide\textquotedblright{} (\emph{kwet}) (KANJIDIC2; Middle Chinese in Baxter's transcription).
\end{cousinweb}

\begin{grammarlens}[title={What you've built}]
One more word for telling what happened.
\end{grammarlens}

\subsection*{Your turn}
\begin{itemize}
\raggedright
  \item \emph{Say it:} \emph{hanketsu}
  \item \emph{Say it:} \emph{hanketsu}, once more
  \item \emph{Say it:} \emph{hanketsu wa raishū desu} — The verdict is next week.
  \item \emph{Recall:} say the Japanese for a trial, then the Japanese for a fine, then say \emph{hanketsu} again
\end{itemize}

\begin{tcolorbox}[breakable,colback=teal!4,colframe=teal!35!black,title={Before you move on}]
What does \ja{はんけつ} mean? (\textquotedblleft{}A verdict; a judgement\textquotedblright{}.) Say it once more.
\end{tcolorbox}
\section[\texorpdfstring{\ja{うったえる} (uttaeru)}{うったえる (uttaeru)}]{\ja{うったえる} (uttaeru) — to sue; to complain (to an authority)}

\label{lesson:JA-C273-uttaeru}



\begin{quote}
Before the new one: say the Japanese for a fine, then the Japanese for a verdict.
\end{quote}

\subsection*{You'll want to know: \ja{うったえる}}
\textbf{\ja{うったえる}} — \emph{uttaeru} — \textquotedblleft{}to sue; to complain (to an authority)\textquotedblright{}.

\begin{quote}
\textbf{\ja{その} \ja{おんなは} \ja{かいしゃを} \ja{うったえました。}} — \emph{sono onna wa kaisha o uttaemashita} — \textquotedblleft{}That woman took the company to court.\textquotedblright{}
\end{quote}

Polite \textbf{\ja{うったえます}}, \emph{uttaemasu}; an ichidan verb. Past (polite): \textbf{\ja{おとこが} \ja{みせを} \ja{うったえました。}}, \emph{otoko ga mise o uttaemashita}.

\begin{cousinweb}
\textbf{\ja{うったえる}} (\emph{uttaeru}) is a native Japanese word (wago), inherited rather than borrowed from Chinese; its kanji was attached later as a kun'yomi (JMdict).
\end{cousinweb}

\begin{grammarlens}[title={What you've built}]
One more word for telling what happened.
\end{grammarlens}

\subsection*{Your turn}
\begin{itemize}
\raggedright
  \item \emph{Say it:} \emph{uttaeru}
  \item \emph{Say it:} \emph{uttaeru}, once more
  \item \emph{Say it:} \emph{sono onna wa kaisha o uttaemashita} — That woman took the company to court.
  \item \emph{Recall:} say the Japanese for a fine, then the Japanese for a verdict, then say \emph{uttaeru} again
\end{itemize}

\begin{tcolorbox}[breakable,colback=teal!4,colframe=teal!35!black,title={Before you move on}]
What does \ja{うったえる} mean? (\textquotedblleft{}To sue; to complain (to an authority)\textquotedblright{}.) Say it once more.
\end{tcolorbox}
\section[\texorpdfstring{\ja{しょうこ} (shōko)}{しょうこ (shōko)}]{\ja{しょうこ} (shōko) — evidence; proof}

\label{lesson:JA-C273-shouko}



\begin{quote}
Before the new one: say the Japanese for a verdict, then the Japanese for to sue.
\end{quote}

\subsection*{You'll want to know: \ja{しょうこ}}
\textbf{\ja{しょうこ}} — \emph{shōko} — \textquotedblleft{}evidence; proof\textquotedblright{}.

\begin{quote}
\textbf{\ja{けいさつは} \ja{しょうこを} \ja{さがしています。}} — \emph{keisatsu wa shōko o sagashite imasu} — \textquotedblleft{}The police are looking for evidence.\textquotedblright{}
\end{quote}

\begin{cousinweb}
\textbf{\ja{しょうこ}} (\emph{shōko}) is Sino-Japanese (kango): two Chinese characters read with their on'yomi, readings borrowed from Chinese. They are \emph{shō}, \textquotedblleft{}evidence\textquotedblright{} (Middle Chinese \emph{tsyingH}) and \emph{ko}, \textquotedblleft{}based on\textquotedblright{} (\emph{kjoH}) (KANJIDIC2; Middle Chinese in Baxter's transcription).
\end{cousinweb}

\begin{grammarlens}[title={What you've built}]
One more word for telling what happened.
\end{grammarlens}

\subsection*{Your turn}
\begin{itemize}
\raggedright
  \item \emph{Say it:} \emph{shōko}
  \item \emph{Say it:} \emph{shōko}, once more
  \item \emph{Say it:} \emph{keisatsu wa shōko o sagashite imasu} — The police are looking for evidence.
  \item \emph{Recall:} say the Japanese for a verdict, then the Japanese for to sue, then say \emph{shōko} again
\end{itemize}

\begin{tcolorbox}[breakable,colback=teal!4,colframe=teal!35!black,title={Before you move on}]
What does \ja{しょうこ} mean? (\textquotedblleft{}Evidence; proof\textquotedblright{}.) Say it once more.
\end{tcolorbox}
\section[(dialogue)]{Looking back: from Keeping Calm to The Trial}

\label{lesson:JA-R273-review-268-273}



\begin{quote}
Say the Japanese for to calm down and a rescue.
\end{quote}

\subsection*{Your turn}
Say these aloud:

\begin{itemize}
\raggedright
  \item to calm down, a rescue, to survive, to help each other, safe
  \item a thief, to steal, last night, to bark, to make a noise
  \item a culprit, an incident, a report to the police, that (one), to touch
  \item to be caught, to chase, finally, to hide (oneself), to catch
  \item a fraud, to trick, to threaten, to hide (something), true
  \item a trial, a fine, a verdict, to sue, evidence
\end{itemize}

\begin{tcolorbox}[breakable,colback=teal!4,colframe=teal!35!black,title={Before you move on}]
To calm down? (\textbf{\ja{おちつく}}, \emph{ochitsuku}.) A rescue? (\textbf{\ja{きゅうじょ}}, \emph{kyūjo}.) To survive? (\textbf{\ja{いきのこる}}, \emph{ikinokoru}.) To help each other? (\textbf{\ja{たすけあう}}, \emph{tasukeau}.) Safe? (\textbf{\ja{あんぜん}}, \emph{anzen}.) A thief? (\textbf{\ja{どろぼう}}, \emph{dorobō}.) To steal? (\textbf{\ja{ぬすむ}}, \emph{nusumu}.) Last night? (\textbf{\ja{ゆうべ}}, \emph{yūbe}.) To bark? (\textbf{\ja{ほえる}}, \emph{hoeru}.) To make a noise? (\textbf{\ja{さわぐ}}, \emph{sawagu}.) A culprit? (\textbf{\ja{はんにん}}, \emph{hannin}.) An incident? (\textbf{\ja{じけん}}, \emph{jiken}.) A report to the police? (\textbf{\ja{つうほう}}, \emph{tsūhō}.) That (one)? (\textbf{\ja{その}}, \emph{sono}.) To touch? (\textbf{\ja{さわる}}, \emph{sawaru}.) To be caught? (\textbf{\ja{つかまる}}, \emph{tsukamaru}.) To chase? (\textbf{\ja{おいかける}}, \emph{oikakeru}.) Finally? (\textbf{\ja{ついに}}, \emph{tsuini}.) To hide (oneself)? (\textbf{\ja{かくれる}}, \emph{kakureru}.) To catch? (\textbf{\ja{つかまえる}}, \emph{tsukamaeru}.) A fraud? (\textbf{\ja{さぎ}}, \emph{sagi}.) To trick? (\textbf{\ja{だます}}, \emph{damasu}.) To threaten? (\textbf{\ja{おどす}}, \emph{odosu}.) To hide (something)? (\textbf{\ja{かくす}}, \emph{kakusu}.) True? (\textbf{\ja{ほんとう}}, \emph{hontō}.) A trial? (\textbf{\ja{さいばん}}, \emph{saiban}.) A fine? (\textbf{\ja{ばっきん}}, \emph{bakkin}.) A verdict? (\textbf{\ja{はんけつ}}, \emph{hanketsu}.) To sue? (\textbf{\ja{うったえる}}, \emph{uttaeru}.) Evidence? (\textbf{\ja{しょうこ}}, \emph{shōko}.)
\end{tcolorbox}
